\documentclass[11pt,a4paper]{article}

% ---------- Packages ----------
\usepackage[margin=1in]{geometry}
\usepackage{amsmath, amssymb, amsthm, mathtools}
\usepackage{microtype}
\usepackage[colorlinks=true,
            linkcolor=blue!50!black,
            urlcolor=blue!50!black,
            citecolor=blue!50!black]{hyperref}

% ---------- Theorem environments ----------
\theoremstyle{plain}
\newtheorem{theorem}{Theorem}[section]
\newtheorem{proposition}[theorem]{Proposition}
\newtheorem{lemma}[theorem]{Lemma}

\theoremstyle{definition}
\newtheorem{definition}[theorem]{Definition}

\theoremstyle{remark}
\newtheorem{remark}[theorem]{Remark}
\newtheorem{example}[theorem]{Example}

% ---------- Math shortcuts ----------
\newcommand{\R}{\mathbb{R}}
\newcommand{\N}{\mathbb{N}}
\newcommand{\E}{\mathbb{E}}
\newcommand{\Var}{\operatorname{Var}}
\newcommand{\F}{\mathcal{F}}
\newcommand{\Prob}{\mathbb{P}}
\newcommand{\Q}{\mathbb{Q}}
\newcommand{\dd}{\mathop{}\!\mathrm{d}}
\newcommand{\eqdef}{\overset{\mathrm{def}}{=}}
\newcommand{\indep}{\perp\!\!\!\perp}
\DeclareMathOperator{\sgn}{sgn}

% ---------- Document metadata ----------
\title{Monte Carlo Foundations:\\
       Estimation, Efficiency, and Random Sampling}
\author{Quant Finance Portfolio --- Phase 2, Block 0}
\date{\today}

\begin{document}
\maketitle

\begin{abstract}
This document opens Phase~2 of the project. It implements no algorithm:
its purpose is to lay the conceptual groundwork on which every Monte Carlo
method developed in the remainder of the phase will rest. Two themes are
developed. First, we formulate derivatives pricing as a numerical
integration problem and analyse the standard Monte Carlo estimator:
strong consistency, asymptotic normality, confidence intervals, the
bias--variance decomposition, and the notion of \emph{computational
efficiency} that motivates the entire programme of variance reduction.
Second, we examine how independent samples are obtained from prescribed
distributions: the desirable properties of pseudorandom number
generators, the inverse-transform method as a universal device, and the
two standard routes to the normal distribution---inversion and
Box--Muller---with attention to the consequences of each choice for the
low-discrepancy methods of Block~3.
\end{abstract}

\tableofcontents

\section{Introduction}
\label{sec:intro}

The pricing of a contingent claim under a chosen risk-neutral measure is
the computation of an expectation:
\begin{equation}
\label{eq:price-as-expectation}
C \;=\; \E^{\Q}\!\left[\, e^{-rT}\,\Phi(S_{[0,T]}) \,\right],
\end{equation}
where $\Phi$ is the discounted payoff functional and $S_{[0,T]}$ denotes
the relevant trajectory of the underlying state. In Phase~1 we treated
the case in which~\eqref{eq:price-as-expectation} admits a closed form,
namely the European call under geometric Brownian motion, where the
expectation reduces to a one-dimensional Gaussian integral evaluated
analytically through the Black--Scholes formula. We obtain greeks by
differentiating that formula, and we recover the implied volatility by
inverting it numerically.

That regime is the exception, not the rule. The functional $\Phi$ may
depend on the entire path (Asian, lookback, barrier options); the
underlying may be high-dimensional (basket options, multi-currency
products); the dynamics may have no tractable transition density (local
volatility, stochastic volatility, jump-diffusion). In all these cases
the expectation in~\eqref{eq:price-as-expectation} must be approximated.
Of the available numerical strategies, Monte Carlo simulation is the
most widely applicable: it requires only the ability to sample paths of
the underlying and evaluate the payoff on those paths.

Phase~2 of this project develops Monte Carlo methods systematically,
from the elementary estimator for European options to least-squares
Monte Carlo for American options. Before any of these algorithms can be
implemented or analysed, two prerequisites must be in place. We must
understand what a Monte Carlo estimator \emph{is}---what it converges
to, how fast, with what uncertainty, and at what cost---and we must
understand how the random samples that drive it are produced. This
document addresses both. Section~\ref{sec:mc-as-integration} treats the
estimator. Section~\ref{sec:sampling} treats the samples.
Section~\ref{sec:roadmap} sketches how the material assembled here
threads through the remaining blocks of the phase.

We presuppose familiarity with the basic apparatus of probability
theory---in particular the law of large numbers, the central limit
theorem, and Slutsky's lemma---and with the risk-neutral pricing
framework as developed, for instance, in Etheridge~\cite{Etheridge}. The
principal reference for the material in this document is
Glasserman~\cite[Chapters~1 and~2]{Glasserman}.

\section{Monte Carlo as Numerical Integration}
\label{sec:mc-as-integration}

\subsection{The pricing problem as an expectation}

We work on a filtered probability space
$(\Omega, \F, \{\F_t\}_{t \in [0,T]}, \Q)$ supporting whatever Brownian
motions and Poisson random measures the model requires. Under the
risk-neutral measure $\Q$, the time-zero price of a European-style
contingent claim with $\F_T$-measurable discounted payoff
$Y \eqdef e^{-rT}\Phi$ is the expectation
\begin{equation}
\label{eq:alpha}
\alpha \;\eqdef\; \E^{\Q}[\,Y\,].
\end{equation}
We will treat $\alpha$ throughout as a deterministic real number---the
quantity to be estimated---and reserve hats for its estimators. The
Phase~1 framework gives us $\alpha$ in closed form for the European call
under Black--Scholes; we recall that $C = S_0\,N(d_1) - K e^{-rT} N(d_2)$
arises precisely as the evaluation of~\eqref{eq:alpha} when $S_T$ is
log-normal under $\Q$.

The Monte Carlo approach to~\eqref{eq:alpha} is structurally trivial: if
we can produce a sequence $Y_1, Y_2, \dots$ of independent copies of
$Y$, we estimate $\alpha$ by their average. The substance of the method
lies entirely in two questions: how good is this estimator, and how do
we produce the samples?

\subsection{The standard Monte Carlo estimator}

Let $Y_1, Y_2, \dots, Y_n$ be independent and identically distributed
copies of $Y$, with finite second moment. The \emph{Monte Carlo
estimator} of $\alpha$ based on $n$ samples is
\begin{equation}
\label{eq:mc-estimator}
\hat{\alpha}_n \;\eqdef\; \frac{1}{n}\sum_{i=1}^{n} Y_i.
\end{equation}

Two properties are immediate. First, $\hat{\alpha}_n$ is unbiased:
$\E[\hat{\alpha}_n] = \alpha$. Second, by the strong law of large
numbers, $\hat{\alpha}_n \to \alpha$ almost surely as $n \to \infty$.
The estimator is therefore strongly consistent. Neither property tells
us how large $n$ must be to attain a given precision; for that we need
the variance of the estimator and its asymptotic distribution.

Writing $\sigma^2 \eqdef \Var(Y) < \infty$, we have, by independence,
\begin{equation}
\label{eq:mc-variance}
\Var(\hat{\alpha}_n) \;=\; \frac{\sigma^2}{n},
\end{equation}
so the standard deviation of the estimator decays like $n^{-1/2}$.

\subsection{Convergence rate and the central limit theorem}
\label{sec:clt}

The central limit theorem upgrades~\eqref{eq:mc-variance} to a
distributional statement:
\begin{equation}
\label{eq:clt}
\sqrt{n}\,(\hat{\alpha}_n - \alpha) \;\xrightarrow{\;d\;}\;
N(0, \sigma^2) \qquad \text{as } n \to \infty.
\end{equation}
This is the analytical backbone of Monte Carlo. Three observations are
in order.

\paragraph{Order of convergence.} The error
$|\hat{\alpha}_n - \alpha|$ is, in a probabilistic sense, of order
$n^{-1/2}$. To halve the error one must \emph{quadruple} the sample
size. This rate is slow when compared with classical quadrature in low
dimensions, but it has a feature that classical quadrature lacks: it
does not deteriorate with dimension.

\paragraph{Dimension-independence.} If $Y = f(U)$ for some
$U \in [0,1]^d$ and a sufficiently nice $f$, the rate
in~\eqref{eq:clt} is $n^{-1/2}$ regardless of $d$. The constant
$\sigma$ does typically depend on $d$, often unfavourably, but the
\emph{order} does not. By contrast, a tensor-product trapezoidal rule
on $[0,1]^d$ achieves error $O(n^{-2/d})$ for smooth integrands, which
becomes useless even in modest dimensions. This is the structural
reason Monte Carlo is the default tool for path-dependent and
multi-asset pricing problems: an Asian option monitored at $252$ dates
is a $252$-dimensional integral, and any product rule is hopeless.

\paragraph{Failure modes.} The CLT requires $\Var(Y) < \infty$. There
are realistic settings where this fails or where it holds but
$\Var(Y)$ is so large that the CLT regime is reached only at sample
sizes exceeding any plausible computational budget. Heavy-tailed payoffs
(e.g.\ deep out-of-the-money digital options under fat-tailed dynamics)
fall into the second category and motivate the importance sampling
techniques of Block~2.

\subsection{Confidence intervals from the sample variance}
\label{sec:ci}

In practice $\sigma$ is not known a priori. The natural plug-in is the
sample standard deviation:
\begin{equation}
\label{eq:sample-variance}
S_n^{\,2} \;\eqdef\; \frac{1}{n-1}\sum_{i=1}^{n}(Y_i - \hat{\alpha}_n)^2,
\qquad
S_n \eqdef \sqrt{S_n^{\,2}}.
\end{equation}
Since $S_n^{\,2} \to \sigma^2$ almost surely, an application of
Slutsky's lemma to~\eqref{eq:clt} yields
\begin{equation}
\label{eq:studentised}
\frac{\sqrt{n}\,(\hat{\alpha}_n - \alpha)}{S_n}
\;\xrightarrow{\;d\;}\; N(0,1).
\end{equation}
For a confidence level $1-\beta$, denoting $z_{\beta/2}$ the upper
$\beta/2$-quantile of the standard normal, the asymptotic
$(1-\beta)$-confidence interval for $\alpha$ is
\begin{equation}
\label{eq:ci}
\hat{\alpha}_n \;\pm\; z_{\beta/2}\,\frac{S_n}{\sqrt{n}}.
\end{equation}
The half-width
\begin{equation}
\label{eq:halfwidth}
h_n \;\eqdef\; z_{\beta/2}\,\frac{S_n}{\sqrt{n}}
\end{equation}
is the standard reportable summary of estimator uncertainty in Monte
Carlo. It is what every implementation in this phase will compute and
return alongside the point estimate.

\begin{remark}
The use of $z_{\beta/2}$ rather than the corresponding $t$-quantile is
purely asymptotic. For $n$ in the thousands, which is the regime of
interest throughout the phase, the difference is negligible.
\end{remark}

\subsection{Bias and the mean square error}
\label{sec:mse}

The estimator~\eqref{eq:mc-estimator} is unbiased by construction, but
this exactness will be lost as soon as we replace the exact dynamics of
the model with a discretisation, or as soon as we estimate a function of
an expectation rather than an expectation itself. Both situations arise
in this phase: the former in Block~1 (Euler--Maruyama and Milstein
schemes), the latter in Block~6 (the nested-expectation structure of
American option pricing).

For a generic estimator $\hat{\alpha}$ of a target $\alpha$, the
mean square error decomposes as
\begin{equation}
\label{eq:mse}
\operatorname{MSE}(\hat{\alpha})
\;=\; \E[(\hat{\alpha} - \alpha)^2]
\;=\; \big(\E[\hat{\alpha}] - \alpha\big)^2 + \Var(\hat{\alpha})
\;=\; \operatorname{Bias}^2(\hat{\alpha}) + \Var(\hat{\alpha}).
\end{equation}
When discretisation enters, $\hat{\alpha}$ depends on a step size
$\Delta t$, and a tradeoff appears: refining $\Delta t$ reduces bias
but, for a fixed computational budget, also reduces the number of
replications and hence increases variance. The MSE-optimal choice of
$\Delta t$ balances the two contributions and will be derived
explicitly in Block~1.

\subsection{Computational efficiency of an estimator}
\label{sec:efficiency}

The remaining piece of conceptual scaffolding is the notion of
\emph{efficiency}, which underpins the entire programme of variance
reduction in Block~2.

Suppose two unbiased estimators $\hat{\alpha}$ and $\hat{\alpha}'$ of
the same quantity $\alpha$ are available, with per-replication
variances $\sigma^2$ and $\sigma'^{\,2}$ and per-replication
computational costs $\tau$ and $\tau'$ respectively. Naively, one
prefers the estimator with smaller variance. This is incorrect: what
matters is the total computational work required to reach a prescribed
precision.

\begin{proposition}[Efficiency comparison]
\label{prop:efficiency}
Let $\hat{\alpha}$ and $\hat{\alpha}'$ be unbiased estimators of
$\alpha$ with per-replication variances $\sigma^2, \sigma'^{\,2}$ and
costs $\tau, \tau'$. Fix a target half-width $\varepsilon$ for the
asymptotic confidence interval~\eqref{eq:ci} at confidence level
$1-\beta$. Then $\hat{\alpha}$ requires less total computational work
than $\hat{\alpha}'$ to attain the prescribed half-width if and only if
\begin{equation}
\label{eq:eff-criterion}
\sigma^2 \tau \;<\; \sigma'^{\,2}\, \tau'.
\end{equation}
\end{proposition}

\begin{proof}
For an estimator with per-replication variance $\sigma^2$, attaining
half-width $\varepsilon$ requires, asymptotically, $n_\varepsilon =
(z_{\beta/2}\,\sigma / \varepsilon)^2$ replications by~\eqref{eq:ci}.
The corresponding total cost is
$n_\varepsilon \cdot \tau = (z_{\beta/2} / \varepsilon)^2 \sigma^2 \tau$.
The factor $(z_{\beta/2}/\varepsilon)^2$ does not depend on the
estimator, so the comparison reduces to that of the products
$\sigma^2 \tau$ and $\sigma'^{\,2}\,\tau'$.
\end{proof}

\begin{definition}
The \emph{computational efficiency} of an unbiased estimator with
variance $\sigma^2$ and per-replication cost $\tau$ is
$(\sigma^2 \tau)^{-1}$.
\end{definition}

\begin{remark}[Why this matters]
Variance reduction techniques almost always increase the per-replication
cost: an antithetic estimator generates one extra path per pair, a
control variate evaluates an additional payoff, importance sampling
involves an extra likelihood ratio, and so on. The relevant question is
not whether the variance went down, but whether the product
$\sigma^2 \tau$ went down. Throughout Block~2 every reported result will
include both factors, and the criterion of success will be a reduction
in the product, not in the variance alone.
\end{remark}

\begin{remark}[Biased estimators]
Proposition~\ref{prop:efficiency} assumes both estimators are unbiased.
For estimators carrying bias~$b$ controlled by some parameter (typically
a discretisation step), the analogous comparison balances MSE against
work, and the MSE-optimal allocation between bias and variance involves
the asymptotic rates of $b$ and $\sigma$ as functions of the parameter.
A representative analysis of this kind is carried out in Block~1 for
the Euler scheme.
\end{remark}

\section{Generating Random Samples}
\label{sec:sampling}

The estimator~\eqref{eq:mc-estimator} requires independent copies of
$Y$. The standard pipeline for producing them is two-stage. First, a
\emph{pseudorandom number generator} produces a deterministic sequence
$u_1, u_2, \dots \in [0,1)$ designed to be statistically
indistinguishable, for the purposes of the simulation, from i.i.d.\
samples of the uniform distribution on $[0,1]$. Second, a
\emph{sampling transformation} maps these uniforms into samples from
the distributions actually required by the model: standard normals for
Brownian increments, exponentials for Poisson interarrivals, and so on.
Both stages admit several constructions, and the choices interact in
ways that matter beyond the present block.

\subsection{Pseudorandom number generators}
\label{sec:prng}

A pseudorandom number generator (PRNG) is a deterministic algorithm
that produces, from an initial \emph{seed}, a long sequence of values in
$[0,1)$ which passes a battery of statistical tests for uniformity and
independence. The output is reproducible: the same seed yields the same
sequence. This determinism is a feature, not a defect: it is what makes
debugging, regression testing, and exact reproduction of numerical
results possible.

The minimum properties demanded of a PRNG suitable for serious
simulation are conventionally:
\begin{enumerate}
\item \textbf{Long period.} The eventual cycle length should exceed by
several orders of magnitude the largest sample size that will be
consumed.
\item \textbf{High-dimensional equidistribution.} For applications that
form $d$-tuples of consecutive outputs and treat them as points in
$[0,1)^d$ (which is exactly what we shall do every time we sample a
multivariate normal), the empirical distribution of these $d$-tuples
should be uniform in $[0,1)^d$ for $d$ as large as is relevant.
\item \textbf{Speed and small state.} Generation must be fast, and the
internal state must be small enough to fit comfortably in cache.
\item \textbf{Statistical quality.} The output must pass standard
empirical test suites (TestU01's BigCrush is the contemporary
benchmark).
\end{enumerate}

\paragraph{What not to use.} The classical \emph{linear congruential
generator} (LCG) defined by $X_{n+1} = (a X_n + c) \bmod m$, normalised
to $u_n = X_n/m$, is instructive as a cautionary example. Marsaglia
showed in 1968 that all $d$-tuples of consecutive outputs of an LCG lie
on at most $(d!\,m)^{1/d}$ parallel hyperplanes in $[0,1)^d$, an
arrangement easily detected by any test sensitive to multidimensional
structure. LCGs with poorly chosen parameters were standard in the
1970s and 1980s and contaminated a great deal of simulation work; some
remain bundled with libraries to this day. They have no place in
financial Monte Carlo.

\paragraph{Acceptable defaults.} Two families of generators are in
widespread use and meet the criteria above:
\begin{itemize}
\item The \emph{Mersenne Twister} (MT19937) of Matsumoto and
Nishimura~\cite{Matsumoto1998} has period $2^{19937}-1$ and is
$623$-distributed to $32$-bit accuracy. It is the default in NumPy's
legacy \texttt{RandomState} and in the C++ standard library
(\texttt{std::mt19937} and \texttt{std::mt19937\_64}).

\item The \emph{permuted congruential generator} family (PCG) of
O'Neill~\cite{ONeill2014} couples an LCG with a non-linear output
permutation, achieving statistical quality comparable to MT19937 with a
much smaller state and faster generation. NumPy's modern
\texttt{Generator} interface uses PCG64 by default.
\end{itemize}
For this project we adopt PCG64 in Python (the NumPy default) and
\texttt{std::mt19937\_64} in C++. The choice is not critical at our
sample sizes; either generator is comfortably above the quality
threshold required.

\paragraph{True randomness.} Hardware-based true random number
generators exist and are appropriate for cryptography, but they are
unsuitable for simulation precisely because their output cannot be
reproduced. A simulation result that depends on a true random sequence
is not falsifiable.

\subsection{The inverse-transform method}
\label{sec:inverse-transform}

Given uniforms, we now turn to producing samples from arbitrary
prescribed distributions. The most general and conceptually cleanest
method is inversion.

\begin{theorem}[Inverse-transform sampling]
\label{thm:inverse-transform}
Let $F: \R \to [0,1]$ be a cumulative distribution function, and define
its generalised inverse
\begin{equation}
F^{-1}(u) \eqdef \inf\{ x \in \R : F(x) \geq u \},
\qquad u \in (0,1).
\end{equation}
If $U \sim U(0,1)$, then $X \eqdef F^{-1}(U)$ has distribution function
$F$.
\end{theorem}

\begin{proof}
By construction, $F^{-1}(u) \leq x$ if and only if $u \leq F(x)$.
Hence $\Prob(X \leq x) = \Prob(F^{-1}(U) \leq x) = \Prob(U \leq F(x))
= F(x)$.
\end{proof}

When $F$ is continuous and strictly increasing, $F^{-1}$ is the
ordinary inverse function. When $F$ has flats (corresponding to gaps in
the support of the distribution) or jumps (corresponding to atoms), the
generalised inverse handles both correctly, and the theorem applies
without modification.

The practical bottleneck of inversion is the evaluation of $F^{-1}$.
For some distributions it is closed-form (the exponential:
$F^{-1}(u) = -\lambda^{-1}\log(1-u)$); for the most important
distribution in finance, the standard normal, it is not, and must be
approximated. Section~\ref{sec:normal-sampling} addresses this.

\paragraph{Why inversion is preferred over alternatives.} Two
properties recommend inversion as the default method.

First, inversion is \emph{coordinate-preserving}: it consumes one
uniform and produces one sample of the target distribution. This means
that if $(U_1, \dots, U_d)$ is a sample from a $d$-dimensional uniform
construction, then $(F^{-1}(U_1), \dots, F^{-1}(U_d))$ is the
corresponding sample of independent variables with distribution $F$,
and the geometric structure of the $d$-dimensional construction is
preserved. This property is essential for quasi-Monte Carlo, the
subject of Block~3, where the points $(U_1, \dots, U_d)$ are drawn from
a low-discrepancy sequence designed to fill $[0,1)^d$ uniformly. Methods
that consume two uniforms per sample (such as Box--Muller, see below)
mix two coordinates of the construction together and destroy this
structure.

Second, inversion is \emph{monotone} in $U$, which makes it compatible
with stratification (Block~2) and with antithetic sampling: the
antithetic of $X = F^{-1}(U)$ is $X' = F^{-1}(1-U)$, which is the
correct antithetic in the sense of being perfectly negatively correlated
with $X$ when $F$ is symmetric.

\subsection{Sampling from the standard normal}
\label{sec:normal-sampling}

Brownian increments are by far the most common random ingredient of
financial Monte Carlo, so the efficient and accurate generation of
standard normal samples deserves dedicated attention.

\subsubsection{Inversion: rational approximations of $\Phi^{-1}$}

The standard normal inverse $\Phi^{-1}$ has no elementary closed form.
What is available, and what every numerical library uses internally, is
a rational approximation to $\Phi^{-1}$ of sufficiently high accuracy
that the residual error is irrelevant compared with Monte Carlo
sampling error. The classical approximation is
Beasley--Springer~\cite{BeasleySpringer1977} as extended by
Moro~\cite{Moro1995}; Wichura's algorithm AS~241~\cite{Wichura1988} is
a widely used alternative; Acklam's approximation~\cite{Acklam2003}
gives full double precision throughout the support and is what
\texttt{scipy.stats.norm.ppf} effectively delivers.

For the implementation work of Phase~2, we will not roll our own:
\texttt{scipy.special.ndtri} (or, equivalently,
\texttt{scipy.stats.norm.ppf}) in Python and
\texttt{boost::math::quantile} on a normal distribution in C++ provide
$\Phi^{-1}$ to machine precision. The conceptual point is that
inversion-based normal sampling is a single call to a transcendental
function: it is a few times slower than the alternative below, and that
cost is the price paid for compatibility with quasi-Monte Carlo.

\subsubsection{Box--Muller and Marsaglia polar}

An elegant alternative converts pairs of uniforms directly into pairs
of independent standard normals, without ever computing $\Phi^{-1}$.

\begin{theorem}[Box--Muller transformation]
\label{thm:box-muller}
Let $U_1, U_2 \sim U(0,1)$ be independent. Define
\begin{align}
\label{eq:box-muller}
X_1 &= \sqrt{-2 \log U_1}\,\cos(2\pi U_2), \\
X_2 &= \sqrt{-2 \log U_1}\,\sin(2\pi U_2).
\end{align}
Then $X_1$ and $X_2$ are independent standard normal random variables.
\end{theorem}

\begin{proof}
Writing $R = \sqrt{X_1^2 + X_2^2}$ and $\Theta = \arctan(X_2/X_1)$,
the joint density of $(R, \Theta)$ that would arise from
$(X_1, X_2) \sim N(0, I_2)$ in polar coordinates is
\[
f_{R,\Theta}(r, \theta) \;=\;
\frac{1}{2\pi}\, r\, e^{-r^2/2},
\qquad r \geq 0,\ \theta \in [0, 2\pi),
\]
which factors as the product of an exponential density on $r^2/2$ and a
uniform density on $[0, 2\pi)$. Equivalently, with $V = R^2/2 \sim
\mathrm{Exp}(1)$ and $\Theta \sim U[0, 2\pi)$ independent,
$(X_1, X_2) = (\sqrt{2V}\cos\Theta, \sqrt{2V}\sin\Theta)$ is standard
bivariate normal. Sampling $V = -\log U_1 \sim \mathrm{Exp}(1)$ by
inversion of the exponential and $\Theta = 2\pi U_2$ by linear scaling
of the uniform recovers~\eqref{eq:box-muller}.
\end{proof}

The Marsaglia polar method is a refinement that replaces the
trigonometric calls with rejection sampling on the unit disc, and is
marginally faster on architectures with slow $\sin$ and $\cos$. The
distinction is irrelevant on modern hardware.

\subsubsection{Trade-offs and the quasi-Monte Carlo consideration}

Box--Muller and inversion produce statistically equivalent samples for
the purposes of classical Monte Carlo. They differ in three respects:
\begin{enumerate}
\item \textbf{Speed.} Box--Muller (or Marsaglia polar) is moderately
faster than inversion, since rational approximation of $\Phi^{-1}$ is
costlier than $\sqrt{}, \log, \cos, \sin$.
\item \textbf{Coordinate use.} Box--Muller consumes two uniforms to
produce two normals. Inversion consumes one uniform to produce one
normal.
\item \textbf{Compatibility with low-discrepancy sequences.} A quasi-MC
construction populates $[0,1)^d$ with deterministic points designed to
have low discrepancy in each coordinate. If we transform these via
inversion coordinatewise, the low-discrepancy property survives. If we
transform them via Box--Muller, we mix coordinates~$1$ and~$2$
nonlinearly through the trigonometric pair, and the equidistribution
property is destroyed in those coordinates. Box--Muller is therefore
incompatible with quasi-Monte Carlo.
\end{enumerate}

For this reason we adopt inversion as the universal sampling strategy
for the entire project. The marginal speed advantage of Box--Muller is
not worth the architectural inflexibility it imposes. Block~3 will
exploit this decision directly.

\subsection{Reproducibility: seeds and streams}
\label{sec:seeds}

Every Monte Carlo result reported in this project must be reproducible
from a recorded seed. This is non-negotiable for two reasons. First,
debugging a wrong result is impossible if the wrong result cannot be
reproduced; an off-by-one error in an estimator has to be examined on a
fixed sequence of paths. Second, regression testing requires that
results not change unexpectedly between commits, which is impossible
without seed control.

In Python, this means consistently constructing generators via
\texttt{rng = numpy.random.default\_rng(seed)} with an explicit integer
\texttt{seed}, and passing \texttt{rng} as a parameter to every
function that consumes randomness. The use of the deprecated
module-level functions \texttt{numpy.random.normal} et al., which
implicitly draw from a global mutable state, is forbidden in this
codebase: it makes thread-safety unattainable and makes the order in
which functions are called silently affect their outputs.

In C++, the analogous discipline is to declare an
\texttt{std::mt19937\_64 rng(seed)} as a local variable in the entry
point of a simulation and pass \texttt{rng} by reference to every
function that consumes randomness. Each function uses the appropriate
\texttt{std::normal\_distribution} or \texttt{std::uniform\_real\_distribution}
object (these are stateless modulo the engine), but the engine itself
is owned by exactly one stack frame and threaded through.

When variance reduction enters in Block~2, certain techniques (in
particular common random numbers, used pervasively in greek estimation)
require that two related estimators consume the \emph{same} sequence of
uniforms. The seed-and-pass-by-reference discipline established now
makes this trivial; an alternative discipline based on global state
makes it impossible.

\section{Roadmap for Phase~2}
\label{sec:roadmap}

The material assembled in this document threads through the remainder
of the phase as follows.

\paragraph{Block~1 (Vanilla Monte Carlo and SDE discretisation).}
Section~\ref{sec:mc-as-integration} provides directly the framework for
the European call MC pricer: an estimator of the
form~\eqref{eq:mc-estimator} applied to discounted payoffs, with the
half-width~\eqref{eq:halfwidth} as the precision report. The
discretisation portion of the block introduces bias for the first time,
and the bias--variance decomposition~\eqref{eq:mse} of
Section~\ref{sec:mse} becomes the analytical tool for choosing time
steps.

\paragraph{Block~2 (Variance reduction).} Every method introduced in
this block---antithetic variates, control variates, stratified
sampling, importance sampling---is justified through the efficiency
criterion of Proposition~\ref{prop:efficiency}: the variance reduction
must dominate the additional per-replication cost. The monotonicity
property of inversion (Section~\ref{sec:inverse-transform}) underlies
the antithetic construction; the seed-control discipline of
Section~\ref{sec:seeds} is required for control variates and for the
common-random-numbers technique used in conjunction with greeks.

\paragraph{Block~3 (Quasi-Monte Carlo).} The decision to adopt inversion
as the universal sampling transformation
(Section~\ref{sec:normal-sampling}) is taken for the express purpose of
this block: low-discrepancy sequences cannot be processed through
Box--Muller without losing their structure. The asymptotic regime
of~\eqref{eq:clt} is replaced by a deterministic Koksma--Hlawka bound,
which improves the order of convergence from $n^{-1/2}$ to nearly
$n^{-1}$ at the cost of dimension dependence in the constant.

\paragraph{Block~4 (Greeks).} Pathwise and likelihood-ratio estimators
of price sensitivities are themselves Monte Carlo estimators, to which
the apparatus of Section~\ref{sec:mc-as-integration} applies directly.
Bumped finite-difference estimators introduce both bias (from the
finite-difference truncation) and variance, and the optimal bump size
balances them according to the same logic as the discretisation step
in Block~1. Common random numbers across the bumped and unbumped
simulations are essential for variance control.

\paragraph{Block~5 (Path-dependent options).} Asian, barrier, and
lookback payoffs are functionals of the entire simulated path. The
machinery of estimator, variance, and confidence interval is unchanged;
the per-replication cost increases, and the variance can be very large
(barrier options especially), making variance reduction techniques from
Block~2 not optional but essential.

\paragraph{Block~6 (American options).} Least-squares Monte Carlo
introduces a nested estimation structure---continuation values are
themselves estimated within a Monte Carlo loop---which produces a bias
that does not vanish as the outer sample size grows. The bias--variance
framework of Section~\ref{sec:mse} is again the relevant analytical
tool, now with a non-trivial bias contribution that must be controlled
by basis selection rather than by sample size alone.

\section{Bibliographic notes}
\label{sec:biblio}

The principal reference throughout Phase~2 is
Glasserman~\cite{Glasserman}; the material of this document corresponds
to Sections~1.1 and Chapter~2 of that text, with a slight reordering
to bring the efficiency criterion into the same section as the basic
estimator analysis. Asmussen and Glynn~\cite{AsmussenGlynn} provide a
more probabilistic and less finance-oriented treatment of the same
material, with a substantially more detailed discussion of estimator
asymptotics and of random number generation; their Chapter~3 is the
recommended supplement for the contents of
Section~\ref{sec:mc-as-integration}, and their Chapter~2 covers the
RNG material of Section~\ref{sec:prng} in greater depth. For PRNGs
specifically, L'Ecuyer~\cite{LEcuyer2012} is a thorough survey from one
of the leading practitioners in the field; the original papers on the
Mersenne Twister~\cite{Matsumoto1998} and the PCG
family~\cite{ONeill2014} are short and readable.

\begin{thebibliography}{99}

\bibitem{Acklam2003}
P.\ J.\ Acklam,
\emph{An algorithm for computing the inverse normal cumulative
distribution function},
Available online (technical note), 2003.

\bibitem{AsmussenGlynn}
S.\ Asmussen and P.\ W.\ Glynn,
\emph{Stochastic Simulation: Algorithms and Analysis},
Springer, 2007.

\bibitem{BeasleySpringer1977}
J.\ D.\ Beasley and S.\ G.\ Springer,
\emph{The percentage points of the normal distribution},
Applied Statistics, 26(1):118--121, 1977.

\bibitem{Etheridge}
A.\ Etheridge,
\emph{A Course in Financial Calculus},
Cambridge University Press, 2002.

\bibitem{Glasserman}
P.\ Glasserman,
\emph{Monte Carlo Methods in Financial Engineering},
Springer, 2004.

\bibitem{LEcuyer2012}
P.\ L'Ecuyer,
\emph{Random Number Generation},
in: J.\ E.\ Gentle, W.\ K.\ H\"ardle, Y.\ Mori (eds.),
Handbook of Computational Statistics, Springer, 2012, pp.\ 35--71.

\bibitem{Matsumoto1998}
M.\ Matsumoto and T.\ Nishimura,
\emph{Mersenne Twister: A 623-dimensionally equidistributed uniform
pseudo-random number generator},
ACM Transactions on Modeling and Computer Simulation, 8(1):3--30, 1998.

\bibitem{Moro1995}
B.\ Moro,
\emph{The full Monte},
Risk, 8(2):57--58, 1995.

\bibitem{ONeill2014}
M.\ E.\ O'Neill,
\emph{PCG: A family of simple fast space-efficient statistically good
algorithms for random number generation},
Technical Report HMC-CS-2014-0905, Harvey Mudd College, 2014.

\bibitem{Wichura1988}
M.\ J.\ Wichura,
\emph{Algorithm AS~241: The percentage points of the normal
distribution},
Applied Statistics, 37(3):477--484, 1988.

\end{thebibliography}

\end{document}
